% !TeX root = ../../Book1.tex
\chapter{多线性代数、外代数与 Jacobi 因子}
\label{b1:app:B}

\begin{chapterabstract}
本附录从多线性映射的线性化出发，构造张量积与外代数，
说明交错型、外向量及其对偶配对之间的关系。
在引入欧氏内积以后，Gram 行列式给出体积，线性映射的外作用给出各维伸缩。
最后区分无向面积与定向体积，为面积公式、余面积公式和流的积分表示补齐代数基础。
\end{chapterabstract}

\begin{learninggoals}
\begin{enumerate}
\item\label{b1:item:app-b-goal-tensor}
从商空间与泛性质两个角度理解张量积，辨认一般张量与一个纯张量的差别。
\item\label{b1:item:app-b-goal-wedge}
完成外代数的构造，掌握外向量、交错型和行列式求值的归一化。
\item\label{b1:item:app-b-goal-gram}
用诱导内积解释 Gram 行列式、非正交坐标的体积因子及可分解条件的作用。
\item\label{b1:item:app-b-goal-jacobian}
区分指定切平面的伸缩与外映射的最大伸缩，统一面积和余面积中的 Jacobi 因子。
\item\label{b1:item:app-b-goal-orientation}
使用顶次外空间固定定向，追踪参数变化与外法向优先约定中的符号。
\end{enumerate}
\end{learninggoals}

\chapref{b1:ch:17}已经用子式写出外积，并由 Gram 行列式计算面积伸缩。
这种坐标方法适合直接进入面积公式，却留下几个代数问题：
为什么换基后仍得到同一个对象？向量的外积与作用于向量的交错型有什么区别？
为何某些文献的公式带有阶乘，而另一些没有？
本附录从这些问题出发，将此前使用的计算放回一个统一构造中。

\ProseParagraphBreak

读者只需熟悉有限维线性代数中的基、商空间、对偶、行列式与实对称矩阵的正交对角化。
前三节的代数构造可以独立阅读；体积的测度解释调用\chapref{b1:ch:17}，
余面积的几何解释调用\chapref{b1:ch:18}。
若只为查阅 Jacobi 因子的约定，可先看\secref{b1:sec:B103}与\secref{b1:sec:B104}；
准备阅读流的定义时，还应掌握交错型的配对和\secref{b1:sec:B105}的定向。
这里始终限于有限维实空间，不讨论无限维张量积的范数与完备化。

\ProseParagraphBreak

黎景辉、白正简、周国晖的《高等线性代数学》\cite{LiBaiZhou2014AdvancedLinear}%
可用于衔接张量积、张量代数与交错型的中文表述。
Simon 的《Lectures on Geometric Measure Theory》\cite[\S25]{Simon1983Lectures}%
则适合核对几何测度论中外向量、交错型与测试形式的使用方式；
本附录补出了该处简述的有限维代数构造。
希望把这些工具继续用于流形上的微分形式与积分，可以阅读杜武亮的%
《An Introduction to Manifolds》\cite{Tu2011Introduction}中微分形式与积分两章。
后一本书在这里是进一步阅读的路径，不承担本附录任何省略的证明。

% !TeX root = ../../Book1.tex
\section{多线性映射与张量积}
\label{b1:sec:B101}

内积、行列式和二阶微分有一个共同特点：把其他变量固定后，它们对剩下的一个变量是线性的。这并不意味着它们对所有变量合在一起线性。例如实数乘法满足 \((2x)(2y)=4xy\)，而不是 \(2xy\)。若仍希望使用核、像、对偶等线性代数工具，就应当改变输入空间，把一对向量所携带的双线性信息组织成一个向量。张量积解决的正是这个问题。

\ProseParagraphBreak

本附录中的向量空间均为有限维实向量空间；代数构造本身不要求内积。对偶 \(V^*\) 表示全部实线性泛函，有限维时它与连续对偶一致。直到\secref{b1:sec:B103}才给空间增加内积，不能在此前用内积把 \(V\) 和 \(V^*\) 默认为同一个对象。

\subsection{把双线性关系变成线性关系}
\label{b1:subsec:B10101}

\begin{definition}\label{b1:def:app-b-multilinear}
设 \(V_1,\ldots,V_k,W\) 为实向量空间。若映射
\[
 B:V_1\times\cdots\times V_k\longrightarrow W
\]
对每个变量分别线性，则称 \(B\) 为\term[duoxianxing-yingshe]{多线性映射}[multilinear map]。当 \(k=2\) 时称为双线性映射；当 \(W=\mathbb R\) 时称为多线性型。
\end{definition}

例如 \(\beta(v,w)=v^{\mathsf T}Aw\) 是双线性型；若给定向量空间 \(V\)，求值映射 \((\lambda,v)\mapsto\lambda(v)\) 是 \(V^*\times V\) 上的双线性型。求值不需要选基，矩阵表达却需要选基。这一区别提示我们：一个合适的构造应当保留双线性映射本身，而不是只保留它在某组坐标中的数组。

\begin{definition}\label{b1:def:app-b-tensor-product}
设 \(V,W\) 为实向量空间。若向量空间 \(P\) 及双线性映射 \(\iota:V\times W\rightarrow P\) 满足以下性质：对任意向量空间 \(Z\) 和双线性映射 \(B:V\times W\rightarrow Z\)，都存在唯一的线性映射 \(\widetilde B:P\rightarrow Z\)，使
\[
 B(v,w)=\widetilde B\bigl(\iota(v,w)\bigr),
\]
则称 \((P,\iota)\) 为 \(V,W\) 的\term[zhangliangji]{张量积}[tensor product]，记为 \(V\otimes W\)，并把 \(\iota(v,w)\) 记为 \(v\otimes w\)。上述唯一分解性质称为这一构造的\term[fanxingzhi]{泛性质}[universal property]。
\end{definition}
\symidx{\(V\otimes W\)：实向量空间的代数张量积}

这里的“对任意 \(Z\)”很重要。它要求同一个空间能够接收所有双线性映射的线性化，而不是为某个预先选好的双线性型临时制造一个表示。张量积的定义并未说每个元素都是一个 \(v\otimes w\)；事实上通常需要有限多个这样的元素相加。

\begin{theorem}[张量积的构造与唯一性]\label{b1:thm:app-b-tensor-construction}
任意有限维实向量空间 \(V,W\) 的张量积存在，且任意两个满足定义的构造之间，存在唯一保持 \(v\otimes w\) 的线性同构。若 \(e_1,\ldots,e_m\) 与 \(f_1,\ldots,f_n\) 分别是 \(V,W\) 的基，则
\[
 \{\,e_i\otimes f_j:1\leq i\leq m,\ 1\leq j\leq n\,\}
\]
是 \(V\otimes W\) 的基。因此 \(\dim(V\otimes W)=mn\)。
\end{theorem}
\begin{proof}
先把 \(V\times W\) 中每一对 \((v,w)\) 当作一个形式基符号 \([v,w]\)，组成实向量空间 \(F\)。这里元素按定义是有限线性组合；即使形式基的指标集合无限，也没有无穷求和问题。令 \(R\) 为下列元素张成的子空间：
\[
 \begin{aligned}
 &[v+v',w]-[v,w]-[v',w],\quad [av,w]-a[v,w],\\
 &[v,w+w']-[v,w]-[v,w'],\quad [v,aw]-a[v,w].
 \end{aligned}
\]
定义 \(P=F/R\)，把 \([v,w]\) 的陪集写成 \(v\otimes w\)。商空间恰好把上述四类关系变为等式，故 \(\iota(v,w)=v\otimes w\) 双线性。

给定 \(B\)，先在线性组合上规定
\[
 \widehat B\left(\sum_{\ell=1}^N a_\ell[v_\ell,w_\ell]\right)
   =\sum_{\ell=1}^N a_\ell \cdot B(v_\ell,w_\ell).
\]
这定义了 \(F\) 上的线性映射。双线性保证每个生成关系都落入它的核，因而 \(R\subseteq\ker\widehat B\)。所以 \(\widehat B\) 唯一下降为 \(P\) 上的线性映射 \(\widetilde B\)。由于所有 \(v\otimes w\) 张成 \(P\)，指定这些元素的像就决定了整个线性映射，得到所需唯一性。

若 \((P',\iota')\) 也满足泛性质，分别把 \(\iota'\) 和 \(\iota\) 作为双线性映射，得到 \(L:P\rightarrow P'\) 与 \(M:P'\rightarrow P\)。复合 \(ML\) 与恒等映射在每个 \(\iota(v,w)\) 上相同；泛性质的唯一性给出 \(ML=I_P\)。另一复合同样为恒等映射，故 \(L\) 为唯一指定的同构。

对基的陈述，写 \(v=\sum_i v_i e_i\)、\(w=\sum_jw_jf_j\)。双线性给出
\[
 v\otimes w=\sum_{i,j}v_iw_j\,e_i\otimes f_j,
\]
所以所列元素张成张量积。令 \(e^i,f^j\) 为对偶基，双线性型
\((v,w)\mapsto e^i(v)f^j(w)\) 诱导线性泛函 \(\ell_{ij}\)。它满足
\(\ell_{ij}(e_a\otimes f_b)=\delta_{ia}\delta_{jb}\)。对任一线性关系逐个作用这些泛函，得到全部系数为零，故所列元素线性无关。若某个原空间是零空间，则双线性已强迫全部生成元为零，维数公式仍然成立。
\end{proof}

取基只是最后验证维数的工具，并未参与 \(F/R\) 的定义。若另一组基产生不同的系数矩阵，这些矩阵仍描述同一个张量。所谓坐标无关，不是拒绝坐标计算，而是说明计算结果所代表的对象并不随坐标改变。

\subsection{一般张量与一个乘积的区别}
\label{b1:subsec:B10102}

形如 \(v\otimes w\) 的张量称为\term[chun-zhangliang]{纯张量}[pure tensor]。一般元素是纯张量的有限和。若 \(v,w\) 都非零，选 \(\lambda\in V^*\)、\(\mu\in W^*\) 使 \(\lambda(v)=\mu(w)=1\)，则由泛性质得到的泛函满足
\[
 (\lambda\otimes\mu)(v\otimes w)=1.
\]
因而非零向量的纯张量不会为零。这里写 \(\lambda\otimes\mu\) 同时预告了对偶张量的自然配对，后面将给出其精确解释。

\ProseParagraphBreak

在 \(V=W=\mathbb R^2\) 中，张量
\[
 \tau=e_1\otimes e_1+e_2\otimes e_2
\]
不是纯张量。否则存在 \(v=(a,b)\)、\(w=(c,d)\)，使
\[
 ac=bd=1,\quad ad=bc=0.
\]
前两个等式迫使 \(a,b,c,d\) 全部非零，与后两个等式矛盾。按所选基把张量排列成矩阵后，纯张量对应秩至多一的矩阵；加法可以增加秩。这也解释了为什么不能把一个张量总当作“两个因子”，更不能未经论证从张量等式中分别比较所谓第一因子和第二因子。

\ProseParagraphBreak

同样，\((v+w)\otimes(v+w)\) 含有两个不同的交叉项
\[
 v\otimes w+w\otimes v.
\]
除非额外取商或施加对称性，不能把它们合并成 \(2v\otimes w\)。下一节的外代数恰好要加入另一种关系：这两个交叉项之和为零，而不是两者相等。

\begin{proposition}\label{b1:prop:app-b-tensor-maps}
若 \(A:V\rightarrow V'\)、\(B:W\rightarrow W'\) 线性，则存在唯一线性映射
\[
 A\otimes B:V\otimes W\longrightarrow V'\otimes W'
\]
满足 \((A\otimes B)(v\otimes w)=Av\otimes Bw\)。它与复合相容：
\[
 (A'\otimes B')(A\otimes B)=(A'A)\otimes(B'B).
\]
此外，自然的交换和结合映射
\[
 \begin{aligned}
 V\otimes W&\longrightarrow W\otimes V,
     &v\otimes w&\longmapsto w\otimes v,\\
 (U\otimes V)\otimes W&\longrightarrow U\otimes(V\otimes W),
     &(u\otimes v)\otimes w&\longmapsto u\otimes(v\otimes w)
 \end{aligned}
\]
都是线性同构。
\end{proposition}
\begin{proof}
第一映射由 \((v,w)\mapsto Av\otimes Bw\) 的双线性和泛性质得到。复合公式两边在每个纯张量上相等，而纯张量张成全空间，因此两边相等。交换映射同样由泛性质得到，交换两次为恒等映射，故可逆。

对结合映射，先固定 \(w\)，把 \((u,v)\mapsto u\otimes(v\otimes w)\) 线性化；所得映射对 \(w\) 仍线性，因为这一性质可在生成元 \(u\otimes v\) 上逐一验证。再对 \((U\otimes V)\times W\) 使用泛性质，便得到所写映射。反向以同样两步定义，两个复合在 \((u\otimes v)\otimes w\) 或 \(u\otimes(v\otimes w)\) 上为恒等，而这些元素分别张成对应空间，所以互为逆映射。
\end{proof}

通过这些指定的同构，可以不写张量积的结合括号，并把
\(V^{\otimes k}\) 称为 \(V\) 的 \(k\) 次张量积，约定 \(V^{\otimes0}=\mathbb R\)。反复应用泛性质，\(V_1\times\cdots\times V_k\) 上的多线性映射恰好对应 \(V_1\otimes\cdots\otimes V_k\) 上的线性映射。这里没有把双线性与联合线性混淆：前者的定义域是直积，后者的定义域是另一个向量空间。

\subsection{张量、线性映射与求值}
\label{b1:subsec:B10103}

\begin{proposition}\label{b1:prop:app-b-tensor-hom}
存在自然线性同构
\[
 V^*\otimes W\longrightarrow\operatorname{Hom}(V,W),
 \quad \lambda\otimes w\longmapsto
      \bigl[v\longmapsto\lambda(v)\cdot w\bigr].
\]
还存在自然同构
\[
 V^*\otimes W^*\longrightarrow (V\otimes W)^*,
 \quad
 (\lambda\otimes\mu)(v\otimes w)=\lambda(v)\cdot\mu(w).
\]
\end{proposition}
\begin{proof}
第一个公式对 \(\lambda,w\) 双线性，故定义了张量积上的线性映射。若 \(e_i\) 是 \(V\) 的基、\(e^i\) 是对偶基，则对每个 \(A:V\rightarrow W\)，张量
\(\sum_i e^i\otimes Ae_i\) 映到 \(A\)，因为
\[
 \sum_i e^i(v)\cdot Ae_i=A\left(\sum_i e^i(v)\cdot e_i\right)=Av.
\]
再选 \(W\) 的基 \(f_j\)，所列基张量 \(e^i\otimes f_j\) 的像正是把第 \(i\) 个基向量送到 \(f_j\)、其余基向量送到零的矩阵单位。矩阵单位线性无关，故映射也单射。

第二个公式先对 \(v,w\) 线性化，再对 \(\lambda,\mu\) 线性化，得到所述映射。它把 \(e^i\otimes f^j\) 送到上一构造中的坐标泛函 \(\ell_{ij}\)，而这些泛函正是基 \(e_i\otimes f_j\) 的对偶基，所以它是同构。两项定义都直接使用求值，不含任意选基；证明中采用的基只用来核验双射性。
\end{proof}

第一项同构把有限秩算子和纯张量联系起来：非零纯张量对应秩一算子，多个纯张量之和对应有限个秩一算子之和。例如 \(V^*\otimes V\) 中的 \(\sum_i e^i\otimes e_i\) 对应 \(I_V\)。这个张量与基无关，因为无论用哪组基，映到的都是同一个恒等映射。

\ProseParagraphBreak

双线性求值还诱导了线性泛函
\[
 c:V^*\otimes V\longrightarrow\mathbb R,\quad
 c(\lambda\otimes v)=\lambda(v).
\]
若 \(A\) 对应于 \(\sum_i e^i\otimes Ae_i\)，则
\[
 c(A)=\sum_i e^i(Ae_i)=\operatorname{tr}A.
\]
因此矩阵的迹可以理解为一个不依赖基的求值运算。这比直接比较两套对角元更接近迹的本质，也为后面把行列式理解成最高次外空间上的作用提供了类比。

\ProseParagraphBreak

本节始终讨论有限维的代数张量积。无限维时，纯张量有限和仍可以构成代数张量积，但赋范空间上还有不同的张量范数及相应完备化，\(\operatorname{Hom}\) 也需要区分全部线性映射和有界线性映射。这里的有限维同构不能不加条件地移用到那些问题。

% !TeX root = ../../Book1.tex
\section{交错型与外代数}
\label{b1:sec:B102}

有向面积不仅随边向量多线性变化，还应当在两条边相同的时候消失。交换两条边会反转方向，给某条边加上另一条边的倍数则不改变面积。行列式已经具有这些性质；外代数把它们从最高维推广到每一个维数，使一条有向线段、一个有向平行四边形以及高维平行多面体能够使用相同的运算规则。

\subsection{交错条件与行列式展开}
\label{b1:subsec:B10201}

\begin{definition}\label{b1:def:app-b-alternating-form}
设 \(k\geq1\)。若多线性型 \(\omega:V^k\rightarrow\mathbb R\) 在任意两个输入向量相同时取值为零，则称 \(\omega\) 为 \(V\) 上的 \(k\) 次\term[jiaocuo-xing]{交错型}[alternating form]。更一般地，对取值于向量空间的多线性映射，也用同一条件定义其交错性。零次交错型约定为实数。
\end{definition}

“交错型”一名可在黎景辉、白正简、周国晖的《高等线性代数学》第四章\cite{LiBaiZhou2014AdvancedLinear}中查到；该书第二、三章分别讨论张量积与张量代数，适合补充阅读本附录的代数构造。

交错条件蕴含交换两个输入会变号。固定其余输入，在所考察的两个位置同时代入 \(v+w\)，由多线性展开得到
\[
 0=\omega(\ldots,v,\ldots,w,\ldots)
   +\omega(\ldots,w,\ldots,v,\ldots).
\]
反过来，若交换两个输入会变号，则将二者取为同一个向量可得 \(2\omega=0\)，在实数域上这就给出交错性。因此任意排列 \(\sigma\) 满足
\[
 \omega(v_{\sigma(1)},\ldots,v_{\sigma(k)})
   =\operatorname{sgn}(\sigma)\omega(v_1,\ldots,v_k).
\]
若输入向量线性相关，把其中一个写成其他向量的线性组合，每项便含有重复输入，故值为零。尤其当 \(k>\dim V\) 时只有零交错型。

设 \(e_1,\ldots,e_n\) 是 \(V\) 的基，\(e^1,\ldots,e^n\) 为对偶基。对递增指标组 \(I=(i_1<\cdots<i_k)\)，定义
\[
 \epsilon^I(v_1,\ldots,v_k)
   =\det\bigl(e^{i_a}(v_b)\bigr)_{a,b=1}^k.
\]
这些都是交错型。把一般输入逐个按基展开，并把不同排列整理到递增次序，就得到
\begin{equation}\label{b1:eq:app-b-alternating-expansion}
 \omega(v_1,\ldots,v_k)
  =\sum_{i_1<\cdots<i_k}
       \omega(e_{i_1},\ldots,e_{i_k})
       \epsilon^I(v_1,\ldots,v_k).
\end{equation}
具体说，重复指标项为零；每一组互不相同的指标产生 \(k!\) 个项，其变号与坐标系数合起来恰是右侧的行列式。由于 \(\epsilon^I(e_{j_1},\ldots,e_{j_k})=\delta_{IJ}\)，这些型线性无关。于是 \(k\) 次交错型空间维数为 \(\binom nk\)，并由其在递增基向量组上的取值唯一决定。

\subsection{由张量代数取商}
\label{b1:subsec:B10202}

将各次张量积组成直和
\[
 T(V)=\bigoplus_{k=0}^{\infty}V^{\otimes k}.
\]
元素只含有限个非零分量。由连接张量因子定义乘法，零次的 \(1\in\mathbb R\) 为单位；结合律来自上一节的结合映射。这个代数称为\term[zhangliang-daishu]{张量代数}[tensor algebra]。它本身并不交换。

现在令 \(I\) 为 \(T(V)\) 中由所有 \(v\otimes v\) 生成的双边理想，也就是所有形如 \(a\otimes v\otimes v\otimes b\) 的元素的有限线性组合，其中 \(a,b\) 可分别属于任意次张量积。因为这些关系都有确定次数，按次数分组后仍是这类关系的和，所以 \(I\) 是各个 \(I\cap V^{\otimes k}\) 的直和。

\begin{definition}\label{b1:def:app-b-exterior-algebra}
上述商代数
\[
 \bigwedge V=T(V)/I
   =\bigoplus_{k=0}^{n}\bigwedge\nolimits^kV,
 \quad
 \bigwedge\nolimits^kV=
   V^{\otimes k}/\bigl(I\cap V^{\otimes k}\bigr)
\]
称为 \(V\) 的\term[waidaishu]{外代数}[exterior algebra]；其中 \(\bigwedge^kV\) 是它的 \(k\) 次分量。商中的乘法记为 \(\wedge\)，称为外积，纯张量的像记为 \(v_1\wedge\cdots\wedge v_k\)。约定 \(\bigwedge^0V=\mathbb R\)。
\end{definition}
\symidx{\(\bigwedge^kV\)：外代数的 \(k\) 次分量}

这里的外积与\subsecref{b1:subsec:170104}的外积是同一个构造；下面的基定理会说明它们的对应关系，故不重复登记“外积”的主索引。由商关系有 \(v\wedge v=0\)，再将 \(v+w\) 代入便得
\[
 v\wedge w=-w\wedge v.
\]
相邻两个因子可以交换变号；因此只要有两个因子相同，就可以将它们移到相邻位置并使整个乘积为零。商关系看似只处理相邻的重复向量，却已蕴含全部交错关系。

\begin{proposition}[外代数的基与交错泛性质]\label{b1:prop:app-b-exterior-basis}
若 \(e_1,\ldots,e_n\) 是 \(V\) 的基，则
\[
 e_I=e_{i_1}\wedge\cdots\wedge e_{i_k},
 \quad i_1<\cdots<i_k,
\]
组成 \(\bigwedge^kV\) 的基。特别地，
\[
 \dim\bigwedge\nolimits^kV=\binom nk,\quad
 \bigwedge\nolimits^1V=V,\quad
 \bigwedge\nolimits^kV=\{0\}\quad(k>n).
\]
每个交错多线性映射 \(B:V^k\rightarrow W\) 唯一对应线性映射
\(\widetilde B:\bigwedge^kV\rightarrow W\)，满足
\(\widetilde B(v_1\wedge\cdots\wedge v_k)=B(v_1,\ldots,v_k)\)。
\end{proposition}
\begin{proof}
把各向量按基展开，所得到的基向量乘积若有重复指标就为零，否则可以通过相邻交换整理成 \(\pm e_I\)。所以这些 \(e_I\) 张成 \(\bigwedge^kV\)。

为了证明无额外关系，对上面定义的 \(\epsilon^I\) 使用张量积的多线性泛性质，得到 \(V^{\otimes k}\) 上的线性泛函。凡是包含相邻的 \(v\otimes v\) 的生成关系，代入后两个输入相同，故它在 \(I\cap V^{\otimes k}\) 上为零。因此该泛函下降到 \(\bigwedge^kV\)，在 \(e_J\) 上的值为 \(\delta_{IJ}\)。逐一测试线性关系便得全部系数为零，证明线性无关。超过 \(n\) 个基向量必有重复，因而此时整个分量为零，也说明定义中的直和确实到 \(n\) 次为止。

最后对 \(B\) 作张量线性化。它同样在所有含相邻重复因子的关系上为零，所以下降为 \(\widetilde B\)。外积 \(v_1\wedge\cdots\wedge v_k\) 张成全空间，故下降后的映射唯一。反方向把任一线性 \(\widetilde B\) 与外积复合，所得映射具有多线性及交错性，因而两种操作互为逆。
\end{proof}

如果 \(v_j=\sum_i t_{ij}e_i\)，刚才的证明还给出坐标公式
\[
 v_1\wedge\cdots\wedge v_k
   =\sum_{|I|=k}\det(T_I)e_I.
\]
这正是\chapref{b1:ch:17}的子式构造。外代数的取商定义说明该公式为何能在任意基下成立；子式公式则提供实际计算的方法。两种表述各自解决一个问题，不需要在“抽象”与“计算”之间二选一。

外代数中形如 \(v_1\wedge\cdots\wedge v_k\) 的元素称为\term[kefenjie-waixiangliang]{可分解外向量}[decomposable multivector]；零元素也包括在内。并非所有外向量都可分解。对 \(\mathbb R^4\) 的标准基，令
\[
 \xi=e_1\wedge e_2+e_3\wedge e_4.
\]
则 \(\xi\wedge\xi=2e_1\wedge e_2\wedge e_3\wedge e_4\ne0\)，而任何 \(v\wedge w\) 与自身的外积都含重复因子，必为零。所以 \(\xi\) 不能表示成一个有向平行四边形。一般外向量更适合理解为这类元素的线性组合。

若 \(\xi\in\bigwedge^kV\)、\(\eta\in\bigwedge^\ell V\)，交换两个块需要交换 \(k\ell\) 对向量，故
\[
 \xi\wedge\eta=(-1)^{k\ell}\eta\wedge\xi.
\]
先对可分解元素证明，再利用双线性推广即可。这条规则没有说所有 \(\xi\wedge\xi\) 都为零；它只在 \(k\) 为奇数时强迫该结论。刚才的二次外向量例子正好展示偶数次数的差别。

\subsection{对偶配对与外积的归一化}
\label{b1:subsec:B10203}

对 \(V^*\) 也作同样的外代数构造。两类外空间的关系必须通过求值说明，不能因为都记作外积，就把向量和作用于向量的泛函混为一谈。

\begin{proposition}\label{b1:prop:app-b-form-pairing}
存在自然同构
\[
 \bigwedge\nolimits^kV^*
   \longrightarrow \bigl(\bigwedge\nolimits^kV\bigr)^*
\]
及其双线性配对，规定为
\[
 \left\langle
  \lambda_1\wedge\cdots\wedge\lambda_k,\,
  v_1\wedge\cdots\wedge v_k
 \right\rangle
 =\det\bigl(\lambda_i(v_j)\bigr)_{i,j=1}^k.
\]
经此同构，\(\bigwedge^kV^*\) 正是 \(V\) 上全部 \(k\) 次交错型的空间。
\end{proposition}
\begin{proof}
右侧对 \(\lambda_i\) 及 \(v_j\) 分别多线性、交错。先对一组变量使用外代数的泛性质，再对另一组使用，便得到良定义的双线性配对。在基
\[
 e^I=e^{i_1}\wedge\cdots\wedge e^{i_k},
 \quad e_J=e_{j_1}\wedge\cdots\wedge e_{j_k}
\]
上，配对值为 \(\delta_{IJ}\)。因此所定义的映射把 \(e^I\) 送到 \(e_I\) 的对偶基，故为同构。交错泛性质或公式\eqref{b1:eq:app-b-alternating-expansion}表明每个交错型都由这些基型的唯一线性组合给出。
\end{proof}
\symidx{\(\bigwedge^kV^*\)：\(V\) 上的 \(k\) 次交错型空间}

以后把 \(\omega(v_1,\ldots,v_k)\) 与
\(\langle\omega,v_1\wedge\cdots\wedge v_k\rangle\) 视为同一数。这里的尖括号是对偶配对；只有引入内积并作指定的识别后，才可进一步将它看作内积。

为了核对不同文献的系数，记 \(S_{k+\ell}\) 为排列群，并用
\(\operatorname{Sh}(k,\ell)\) 表示保持前 \(k\) 个位置及后 \(\ell\) 个位置各自顺序的排列。若 \(\alpha,\beta\) 的次数分别为 \(k,\ell\)，则由上述约定得到
\begin{equation}\label{b1:eq:app-b-wedge-shuffle}
 \begin{aligned}
 (\alpha\wedge\beta)(v_1,\ldots,v_{k+\ell})
  =\sum_{\sigma\in\operatorname{Sh}(k,\ell)}
     &\operatorname{sgn}(\sigma)\,
       \alpha(v_{\sigma(1)},\ldots,v_{\sigma(k)})\\
     &{}\cdot
       \beta(v_{\sigma(k+1)},\ldots,v_{\sigma(k+\ell)}).
 \end{aligned}
\end{equation}
证明只需先取 \(\alpha,\beta\) 为一形式的外积，将大行列式按前 \(k\) 行展开；每个被选中的列组对应一个保序排列，两个小行列式分别给出 \(\alpha,\beta\) 的值。一般情形再按外空间的基线性展开。等价地，可以对全部排列求和后除以 \(k!\ell!\)：同一保序排列对应的内部排列恰有 \(k!\ell!\) 个，内部变号分别被两个交错型吸收。

例如
\[
 (\lambda\wedge\mu)(v,w)
   =\lambda(v)\mu(w)-\lambda(w)\mu(v),
\]
没有额外的 \(\frac12\)。若把一般 \(r\) 次多线性型 \(T\) 的交错化定义为
\[
 \operatorname{Alt}T(v_1,\ldots,v_r)
  =\frac1{r!}\sum_{\sigma\in S_r}
       \operatorname{sgn}(\sigma)
       T(v_{\sigma(1)},\ldots,v_{\sigma(r)}),
\]
则本书的外积满足
\[
 \alpha\wedge\beta
   =\frac{(k+\ell)!}{k!\ell!}\,
      \operatorname{Alt}(\alpha\otimes\beta).
\]
这个因子来自两种不同的安排：\(\operatorname{Alt}\) 是平均，而外积要使基型与基外向量的配对恰为 \(1\)。以后积分微分形式时沿用后一约定。

类似地，把 \(\bigwedge^kV\) 识别为张量空间中的交错部分时，可送
\[
 v_1\wedge\cdots\wedge v_k
 \longmapsto
 \frac1{k!}\sum_{\sigma\in S_k}
   \operatorname{sgn}(\sigma)
   v_{\sigma(1)}\otimes\cdots\otimes v_{\sigma(k)}.
\]
在外代数商映射之后，这个表达又回到原来的外积，因此确为单射。这个识别中的 \(1/k!\) 与型的行列式求值公式并不冲突，但提醒我们：不能既使用张量子空间的平均定义，又未经换算套用另一种外积归一化。

% !TeX root = ../../Book1.tex
\section{诱导内积、Gram 行列式与体积}
\label{b1:sec:B103}

到目前为止，外积只记录代数上的方向和变号规则。它没有长度，因而也没有面积。现在给 \(V\) 选定欧氏内积 \(\langle\cdot,\cdot\rangle_V\)，并记向量长度为 \(|v|\)。本节把这个内积传到各次外空间；传递规则必须使一组标准正交向量张成的单位平行多面体仍有单位体积，而不能因外积归一化而多出一个阶乘。

\subsection{外空间上的欧氏内积}
\label{b1:subsec:B10301}

\begin{theorem}[外空间的诱导内积]\label{b1:thm:app-b-gram-inner-product}
存在唯一的内积，使任意标准正交基 \(e_1,\ldots,e_n\) 所给出的 \(e_I\) 都成为 \(\bigwedge^kV\) 的标准正交基。对任意 \(v_1,\ldots,v_k,w_1,\ldots,w_k\in V\)，有
\begin{equation}\label{b1:eq:app-b-exterior-inner-product}
 \left\langle
    v_1\wedge\cdots\wedge v_k,\,
    w_1\wedge\cdots\wedge w_k
 \right\rangle
  =\det\bigl(\langle v_i,w_j\rangle_V\bigr)_{i,j=1}^k.
\end{equation}
特别地，
\[
 |v_1\wedge\cdots\wedge v_k|
  =\sqrt{\det\bigl(\langle v_i,v_j\rangle_V\bigr)_{i,j=1}^k}.
\]
零次空间 \(\bigwedge^0V=\mathbb R\) 取通常的实内积。
\end{theorem}
\begin{proof}
先选定一组标准正交基，声明 \(e_I\) 标准正交，就得到一个正定内积。若 \(A,B\) 分别是 \(v_i,w_j\) 的列坐标矩阵，则外积的子式展开给出
\[
 \left\langle\sum_I\det(A_I)e_I,\sum_J\det(B_J)e_J\right\rangle
   =\sum_I\det(A_I)\det(B_I)
   =\det(A^{\mathsf T}B).
\]
最后一个等式是\subsecref{b1:subsec:170104}证明的 Cauchy--Binet 公式；右侧矩阵的第 \(i,j\) 项正是 \(\langle v_i,w_j\rangle_V\)，从而得到所述公式。

公式右侧与最初所选的标准正交基无关，而所有可分解外向量张成 \(\bigwedge^kV\)。所以该公式唯一决定整个双线性型，也证明换用另一组标准正交基会得到同一内积。将 \(w_i=v_i\) 并取非负平方根，便得到范数公式。零次空间的空行列式取 \(1\)，与约定一致。
\end{proof}

Simon 的《Lectures on Geometric Measure Theory》\cite[\S25, pp. 129--131]{Simon1983Lectures}在介绍流之前采用了这一标准正交基约定；本节的证明补出由外积坐标到内积不变性的代数步骤。特别是，这里的 \(|\xi|\) 是外空间作为有限维欧氏空间的长度，不预先赋予“一般外向量所携带的总面积”这一含义。

在上一节给出的张量交错部分识别下，情况略有不同。如果张量空间中各个 \(e_{i_1}\otimes\cdots\otimes e_{i_k}\) 标准正交，那么
\[
 \left|
   \frac1{k!}\sum_{\sigma\in S_k}
       \operatorname{sgn}(\sigma)
       e_{\sigma(1)}\otimes\cdots\otimes e_{\sigma(k)}
 \right|^2=\frac1{k!}.
\]
它的外代数对应物 \(e_1\wedge\cdots\wedge e_k\) 的长度却是 \(1\)。两种空间之间的代数同构不自动是等距同构。若从张量子空间出发定义外空间内积，就必须作相应的归一化。

\begin{proposition}\label{b1:prop:app-b-decomposable-plane}
对 \(1\leq k\leq n\)，外积 \(v_1\wedge\cdots\wedge v_k\) 非零当且仅当各 \(v_i\) 线性无关。若
\(\xi=v_1\wedge\cdots\wedge v_k\ne0\)，则它唯一确定平面
\[
 P_\xi=\operatorname{span}\{v_1,\ldots,v_k\}
      =\{\,v\in V:v\wedge\xi=0\,\}.
\]
存在该平面的一组标准正交基 \(q_1,\ldots,q_k\)，使
\[
 \xi=|\xi|\,q_1\wedge\cdots\wedge q_k.
\]
\end{proposition}
\begin{proof}
若各向量线性相关，交错性使外积为零；若线性无关，将它们扩充成 \(V\) 的一组基，则其外积是该组基对应的一个外空间基向量，所以非零。

把 \(v_1,\ldots,v_k\) 扩充为基。属于其张成空间的向量与 \(\xi\) 外积为零；不属于这一空间的向量至少有一个补空间方向的非零系数，与 \(\xi\) 外积后得到互不相同的基外向量的非平凡线性组合，因而非零。这证明了平面的内在刻画。

在这个平面内选标准正交基 \(q_i\)，将 \(v_j=\sum_i a_{ij}q_i\) 代入外积，得到
\[
 \xi=(\det A)\,q_1\wedge\cdots\wedge q_k.
\]
由诱导范数，\(|\xi|=|\det A|\)。若行列式为负，就把 \(q_1\) 改为 \(-q_1\)，便得到正的系数。这个符号选择将在定向的定义中发挥作用。
\end{proof}

长度相同并不意味着代表同一个平面；同一个平面也可以有相反方向。例如 \(e_1\wedge e_2\) 与 \(-e_1\wedge e_2\) 张成同一二维平面、具有相同长度，却是不同外向量。外积因此同时保存了一个平面、一个尺度和一个方向，取绝对长度时才丢弃后两者中的方向信息。

\subsection{体积与非正交坐标}
\label{b1:subsec:B10302}

设 \(v_1,\ldots,v_k\) 线性无关，考虑其生成的平行多面体
\[
 Q=\left\{\sum_{j=1}^k t_jv_j:0\leq t_j\leq1\right\}.
\]
在平面 \(P=\operatorname{span}\{v_j\}\) 上使用由环境内积诱导的距离，并采用\chapref{b1:ch:15}的 Hausdorff 测度归一化。由线性面积公式，
\begin{equation}\label{b1:eq:app-b-parallelepiped-volume}
 \mathcal H^k(Q)
 =|v_1\wedge\cdots\wedge v_k|
 =\sqrt{\det\bigl(\langle v_i,v_j\rangle_V\bigr)}.
\end{equation}
若向量相关，则 \(Q\) 落在低维平面中，两侧都为零。这里先有外代数内积，再借\chapref{b1:ch:17}已证的测度公式解释其长度，因而没有用直观体积来反过来证明代数公式。

还可直接看到“底乘高”的结构。将 \(v_j\) 依次分解为其在前面向量张成空间中的投影与正交余量 \(w_j\)。在外积中，投影项含有重复方向而消失，所以
\[
 v_1\wedge\cdots\wedge v_k=w_1\wedge\cdots\wedge w_k,
 \quad
 |v_1\wedge\cdots\wedge v_k|=\prod_{j=1}^k|w_j|.
\]
每个 \(|w_j|\) 都是新边相对于已有底面的垂直高度。由 \(|w_j|\leq|v_j|\) 得到
\[
 |v_1\wedge\cdots\wedge v_k|\leq\prod_{j=1}^k|v_j|.
\]
若每个 \(v_j\ne0\)，且它们线性无关，那么等号成立恰好要求每一步投影为零，也就是各边两两正交；线性相关时左侧为零，若右侧非零便不可能取等。

在非标准正交基 \(b_1,\ldots,b_n\) 下，记
\[
 G=(\langle b_i,b_j\rangle_V)_{i,j=1}^n.
\]
若 \(v_j\) 的列坐标组成 \(n\times k\) 矩阵 \(A\)，则
\[
 |v_1\wedge\cdots\wedge v_k|^2=\det(A^{\mathsf T}GA).
\]
坐标中的 \(A^{\mathsf T}A\) 只有在 \(G=I_n\) 时才是正确的 Gram 矩阵。尤其对全维平行多面体，
\[
 \mathcal H^n(Q)=\sqrt{\det G}\,|\det A|.
\]
这不是新的测度归一化，而是坐标基本身所围出的单位盒子已经具有体积 \(\sqrt{\det G}\)。

例如在 \(\mathbb R^2\) 中取 \(b_1=(1,0)\)、\(b_2=(1,1)\)。向量 \(b_2\) 的坐标为 \((0,1)\)，但真实长度为 \(\sqrt2\)；由于
\[
 G=\begin{pmatrix}1&1\\1&2\end{pmatrix},
\]
正确的坐标长度是 \(\sqrt{x^{\mathsf T}Gx}\)，而不是坐标数组的通常长度。另一方面 \(\det G=1\)，这组基围出的平行四边形面积恰为 \(1\)。长度失真与总体积失真是不同问题。

\subsection{外积的估计与一般外向量}
\label{b1:subsec:B10303}

\begin{proposition}\label{b1:prop:app-b-wedge-estimates}
若 \(\xi\in\bigwedge^kV\)、\(\eta\in\bigwedge^\ell V\)，且 \(k+\ell\leq n\)，则
\[
 |\xi\wedge\eta|
 \leq\sqrt{\binom{k+\ell}{k}}\,
        |\xi|\cdot|\eta|.
\]
若 \(\xi\) 可分解，则有更强的估计
\[
 |\xi\wedge\eta|\leq|\xi|\cdot|\eta|.
\]
当 \(k+\ell>n\) 时，外积为零。
\end{proposition}
\begin{proof}
选标准正交基，写 \(\xi=\sum_I\xi_Ie_I\)、\(\eta=\sum_J\eta_Je_J\)。对每个大小为 \(k+\ell\) 的指标组 \(K\)，外积的 \(e_K\) 系数是
\[
 \sum_{\substack{I\cup J=K\\I\cap J=\varnothing}}
       s_{I,J}\xi_I\eta_J,\quad s_{I,J}\in\{-1,1\}.
\]
和式有 \(\binom{k+\ell}{k}\) 项。对这个有限和用 Cauchy--Schwarz，再对 \(K\) 求和，得到
\[
 |\xi\wedge\eta|^2
 \leq\binom{k+\ell}{k}
      \sum_{I\cap J=\varnothing}\xi_I^2\eta_J^2
 \leq\binom{k+\ell}{k}
      \left(\sum_I\xi_I^2\right)
      \left(\sum_J\eta_J^2\right).
\]

若 \(\xi\ne0\) 可分解，由上一命题可选择标准正交基使
\(\xi=|\xi|e_1\wedge\cdots\wedge e_k\)。展开 \(\eta\) 后，凡含有前 \(k\) 个指标的项都被外积消去，其余项仍对应不同的标准正交基外向量。因此
\[
 |\xi\wedge\eta|^2
   =|\xi|^2
      \sum_{J\cap\{1,\ldots,k\}=\varnothing}\eta_J^2
   \leq|\xi|^2\cdot|\eta|^2.
\]
零元素情形直接成立；次数大于维数时则由外空间为零得到结论。
\end{proof}

可分解条件在第二个估计中有实际作用。例如在 \(\mathbb R^6\) 中取
\[
 \xi=e_1\wedge e_2+e_3\wedge e_4+e_5\wedge e_6.
\]
有 \(|\xi|=\sqrt3\)，而
\[
 \xi\wedge\xi
  =2(e_1\wedge e_2\wedge e_3\wedge e_4
     +e_1\wedge e_2\wedge e_5\wedge e_6
     +e_3\wedge e_4\wedge e_5\wedge e_6)
\]
的长度为 \(2\sqrt3>3=|\xi|^2\)。所以外空间的欧氏长度不能不分对象地当成一个使乘法常数总为 \(1\) 的范数。

同样，在 \(\mathbb R^4\) 中，\(\xi=e_1\wedge e_2+e_3\wedge e_4\) 的欧氏长度为 \(\sqrt2\)，却不能据此说“两张互相正交的单位面片总面积为 \(\sqrt2\)”。这个代数和只保留一个外向量；几何上的两块面片若作为不同的积分对象相加，还必须说明支集、重数以及怎样测试它们。\appref{b1:app:D}将据此引入流的质量和余质量；在可分解外向量上有关长度相容，但一般外向量上的质量范数与本节欧氏范数不得混用。

% !TeX root = ../../Book1.tex
\section{线性映射的外作用与 Jacobi 因子}
\label{b1:sec:B104}

在普通线性代数中，映射 \(A:V\rightarrow W\) 逐个作用于向量。若我们关心的对象是一组向量共同张成的面积元，就应让 \(A\) 同时作用于它们，再比较外积的长度。本节从这一简单操作出发，统一\chapref{b1:ch:17}的面积因子和\chapref{b1:ch:18}的余面积因子，同时说明中间维数的最大伸缩与某个指定切平面的伸缩为何不是同一个数。

\subsection{诱导映射、子式与拉回}
\label{b1:subsec:B10401}

\begin{definition}\label{b1:def:app-b-exterior-map}
给定线性映射 \(A:V\rightarrow W\)，由
\[
 \bigwedge\nolimits^k A:
     \bigwedge\nolimits^kV\longrightarrow\bigwedge\nolimits^kW,
 \quad
 v_1\wedge\cdots\wedge v_k
       \longmapsto Av_1\wedge\cdots\wedge Av_k
\]
定义它在 \(k\) 次外空间上的诱导映射。零次诱导映射规定为 \(\mathbb R\) 上的恒等映射。
\end{definition}
\symidx{\(\bigwedge^kA\)：线性映射在外空间上的诱导映射}

这一定义的良定义性直接来自交错泛性质：右侧对输入多线性，而且任意两个输入相同时为零。它不依赖一个外向量被写成哪些可分解外向量之和。对可复合映射 \(A,B\)，有
\[
 \bigwedge\nolimits^k(BA)
  =\bigl(\bigwedge\nolimits^kB\bigr)
       \bigl(\bigwedge\nolimits^kA\bigr),\quad
 \bigwedge\nolimits^kI=I.
\]
两式都可以先在生成元上验证，再用线性性推广。若 \(A\) 可逆，诱导映射也可逆，其逆就是 \(\bigwedge^k(A^{-1})\)。

在 \(V,W\) 的基下，把 \(A\) 的矩阵记为 \(M\)。则 \(\bigwedge^kA\) 的行、列分别按大小为 \(k\) 的递增指标组 \(I,J\) 编号，其矩阵元为
\[
 \bigl[\bigwedge\nolimits^kA\bigr]_{I,J}=\det(M_{I,J}).
\]
这是把 \(Ae_{j_1}\wedge\cdots\wedge Ae_{j_k}\) 按目标外空间基展开的直接结果。复合公式因此包含所有阶数子式的 Cauchy--Binet 关系，而不只是最高阶行列式的乘法公式。

交错型的运算方向相反。若 \(\omega\in\bigwedge^kW^*\)，规定
\[
 (A^*\omega)(v_1,\ldots,v_k)
      =\omega(Av_1,\ldots,Av_k).
\]
这个交错型称为 \(\omega\) 沿 \(A\) 的\term[lahui]{拉回}[pullback]。它把目标空间上用来测试面积元的型变为定义域上的型。因此 \((BA)^*=A^*B^*\)，次序与向量的传递相反。

\begin{proposition}\label{b1:prop:app-b-exterior-adjoint}
上述拉回满足
\[
 \langle A^*\omega,\xi\rangle
    =\langle\omega,(\bigwedge\nolimits^kA)\xi\rangle,
 \quad
 A^*(\alpha\wedge\beta)=A^*\alpha\wedge A^*\beta.
\]
若 \(V,W\) 具有欧氏内积，用 \(A^\dagger:W\rightarrow V\) 表示欧氏伴随，则
\[
 \bigl(\bigwedge\nolimits^kA\bigr)^\dagger
      =\bigwedge\nolimits^k(A^\dagger).
\]
\end{proposition}
\begin{proof}
第一个等式在 \(\xi=v_1\wedge\cdots\wedge v_k\) 上就是拉回的定义，故对一般 \(\xi\) 由线性性成立。对一形式的外积，由行列式配对得到
\[
 A^*(\lambda_1\wedge\cdots\wedge\lambda_k)
    =(\lambda_1\circ A)\wedge\cdots\wedge(\lambda_k\circ A).
\]
可分解型张成全部型空间，再将 \(\alpha,\beta\) 展开，便证明与外积的相容性。

最后取 \(\xi=v_1\wedge\cdots\wedge v_k\)、
\(\eta=w_1\wedge\cdots\wedge w_k\)，由诱导内积公式，
\[
 \begin{aligned}
 \langle(\bigwedge\nolimits^kA)\xi,\eta\rangle
  &=\det\bigl(\langle Av_i,w_j\rangle_W\bigr)\\
  &=\det\bigl(\langle v_i,A^\dagger w_j\rangle_V\bigr)\\
  &=\langle\xi,(\bigwedge\nolimits^kA^\dagger)\eta\rangle.
 \end{aligned}
\]
再用双线性将等式推广到任意外向量，按伴随的唯一性即得结论。
\end{proof}

记号 \(A^*\) 在本节表示对型的拉回；\(A^\dagger\) 表示由两个内积确定的向量映射。通过欧氏 Riesz 对应识别向量与一形式后，两者相互对应，但在没有内积时只有前一种构造。用标准正交坐标表示 \(A^\dagger\)，其矩阵就是熟悉的 \(A^{\mathsf T}\)。

\subsection{奇异值控制的多维伸缩}
\label{b1:subsec:B10402}

\begin{theorem}[外映射的奇异值]\label{b1:thm:app-b-exterior-singular-values}
设 \(\dim V=n\)、\(\dim W=m\)，\(A:V\rightarrow W\) 为欧氏空间之间的线性映射，非零奇异值按
\(\sigma_1\geq\cdots\geq\sigma_r>0\) 排列，其中 \(r=\operatorname{rank}A\)。当 \(1\leq k\leq r\) 时，\(\bigwedge^kA\) 的非零奇异值为
\[
 \sigma_{i_1}\cdots\sigma_{i_k},
 \quad 1\leq i_1<\cdots<i_k\leq r,
\]
并且
\[
 \|\bigwedge\nolimits^kA\|
     =\sigma_1\cdots\sigma_k.
\]
当 \(k>r\) 时，\(\bigwedge^kA=0\)。
\end{theorem}
\begin{proof}
对自伴非负算子 \(A^\dagger A\) 作有限维正交对角化，取标准正交特征基 \(e_i\)，将正特征值写为 \(\sigma_i^2\)。对 \(i\leq r\)，令 \(f_i=Ae_i/\sigma_i\)，则
\[
 \langle f_i,f_j\rangle_W
 =\frac{\langle e_i,A^\dagger Ae_j\rangle_V}{\sigma_i\sigma_j}
 =\delta_{ij}.
\]
把 \(f_1,\ldots,f_r\) 扩充为 \(W\) 的标准正交基。对 \(i>r\)，由
\(|Ae_i|^2=\langle e_i,A^\dagger Ae_i\rangle_V=0\) 得 \(Ae_i=0\)。

由外映射的定义，对 \(I=(i_1<\cdots<i_k)\) 有
\[
 (\bigwedge\nolimits^kA)e_I
  =\begin{cases}
     (\sigma_{i_1}\cdots\sigma_{i_k})f_I,&i_k\leq r,\\
     0,&i_k>r.
   \end{cases}
\]
源与目标的这些外基都标准正交，所以诱导映射已经被写成对角的长方形矩阵。非零对角元就是非零奇异值。对单位外向量，各坐标平方之和为 \(1\)，像的长度不超过最大的对角元；取 \(e_1\wedge\cdots\wedge e_k\) 恰好达到它，从而得到算子范数。若 \(k>r\)，每个基外向量都被送到零。
\end{proof}

这里的最大值在一个可分解单位外向量上达到。因此算子范数虽然允许在所有单位外向量上取上确界，结果仍可只用单位平行多面体求得：
\[
 \|\bigwedge\nolimits^kA\|
  =\max_{\substack{v_1,\ldots,v_k\\\text{标准正交}}}
       |Av_1\wedge\cdots\wedge Av_k|.
\]
不过，某个指定 \(k\) 维平面上的伸缩未必达到最大值。例如
\(A=\operatorname{diag}(3,2,1)\) 在 \(e_1,e_2\) 张成的平面上把面积放大 \(6\) 倍，在 \(e_2,e_3\) 张成的平面上只放大 \(2\) 倍。不能把两种伸缩都写成同一个“二维行列式”。

对指定 \(k\) 维平面 \(P\subseteq V\)，选择它的一组标准正交基 \(q_i\)，记
\[
 J_k(A;P)=|Aq_1\wedge\cdots\wedge Aq_k|.
\]
若换用另一组标准正交基，外积只乘以 \(\pm1\)，所以这个数与选择无关；它就是 \(A|_P\) 的面积因子。若 \(A|_P\) 单射，令 \(Q=A(P)\)，则由外映射的复合和归一化可分解外向量，有
\[
 J_k(BA;P)=J_k(B;Q)J_k(A;P).
\]
若 \(A|_P\) 秩不足，两边的面积都为零，但此时 \(Q\) 不再是 \(k\) 维平面，不能继续把 \(J_k(B;Q)\) 当作上面定义的量。一般的不等式
\[
 J_k(BA;P)
 \leq\|\bigwedge\nolimits^kB\|\,J_k(A;P)
 \leq\|B\|^kJ_k(A;P)
\]
则不需要单射假设。

\subsection{面积因子与余面积因子的统一}
\label{b1:subsec:B10403}

先设 \(\dim V=k\leq\dim W\)，并取 \(V\) 的单位顶次外向量 \(\tau=e_1\wedge\cdots\wedge e_k\)。因为 \(\bigwedge^kV\) 是一维空间，
\[
 J_k(A)=|(\bigwedge\nolimits^kA)\tau|
       =\|\bigwedge\nolimits^kA\|
       =\sqrt{\det(A^\dagger A)}.
\]
在标准正交坐标下，这就是\chapref{b1:ch:17}的列 Gram 公式
\(\sqrt{\det(A^{\mathsf T}A)}\)。若 \(A\) 单射，诱导外映射把定义域的单位面积元变成像平面的面积元；取长度就给出了无向面积伸缩。

若改为 \(\dim W=k\leq\dim V\)，取 \(W\) 的单位顶次外向量 \(\zeta\)。此时自然的量为
\[
 |(\bigwedge\nolimits^kA^\dagger)\zeta|
   =\sqrt{\det(AA^\dagger)}
   =\|\bigwedge\nolimits^kA\|.
\]
在标准正交坐标下，它是\chapref{b1:ch:18}的行 Gram 公式
\(\sqrt{\det(AA^{\mathsf T})}\)。这不是把 \(A^\dagger A\) 的全部特征值相乘；当定义域维数更大时，该算子必有零特征值。我们只取与目标维数相应的非零体积伸缩。

当 \(A\) 满射时，记 \(N=(\ker A)^\perp\)。限制
\(A|_N:N\rightarrow W\) 是 \(k\) 维空间之间的同构，而 \(A\) 在核方向上为零。取核及其正交补的标准正交基后，矩阵写成 \([0\ C]\)，故
\[
 \sqrt{\det(AA^\dagger)}=|\det C|=J_k(A;N).
\]
因此余面积因子计量法方向上的伸缩，核方向的体积则由水平集内的 Hausdorff 测度计算。\chapref{b1:ch:18}的线性余面积公式正是把这两个方向分开积分。

例如 \(A:\mathbb R^3\rightarrow\mathbb R^2\) 定义为
\[
 A(x,y,z)=(x+z,y).
\]
它的行 Gram 矩阵为 \(\operatorname{diag}(2,1)\)，所以 \(J_2(A)=\sqrt2\)。但在水平平面 \(\operatorname{span}(e_1,e_2)\) 上，限制映射的面积因子为 \(1\)。达到最大伸缩的平面是核的正交补，它由
\((e_1+e_3)/\sqrt2,e_2\) 张成。余面积因子不是任意一个横截面上的面积因子；选择非正交截面还会带来相应的几何修正。

对一般坐标，设定义域维数为 \(k\)，定义域基的 Gram 矩阵为 \(G\)，目标基的 Gram 矩阵为 \(H\)，而 \(A\) 的坐标矩阵为 \(M\)。定义域基外积的长度为 \(\sqrt{\det G}\)，其像的长度为 \(\sqrt{\det(M^{\mathsf T}HM)}\)，所以
\begin{equation}\label{b1:eq:app-b-coordinate-jacobian}
 J_k(A)
   =\sqrt{\frac{\det(M^{\mathsf T}HM)}{\det G}}.
\end{equation}
这就是列 Gram 公式的坐标无关内容。在任意基下只写 \(\sqrt{\det(M^{\mathsf T}M)}\)，会把坐标数组当作标准正交坐标；遗漏的并非高阶误差，而是定义域与目标域的实际度量。

最后，若 \(f:U\subseteq\mathbb R^k\rightarrow\mathbb R^n\) 在 \(x\) 可微，其微分 \(Df(x)\) 是线性映射，故
\[
 J_kf(x)
 =|\partial_1f(x)\wedge\cdots\wedge\partial_kf(x)|.
\]
若再作可微参数变换 \(x=\phi(s)\)，普通链式法则与外映射的复合给出
\[
 J_k(f\circ\phi)(s)
   =J_kf\bigl(\phi(s)\bigr)|\det D\phi(s)|.
\]
此处 \(\phi\) 在两端都是 \(k\) 维参数空间，故顶次外作用恰为乘以一个行列式，秩不足时两边仍同时为零。这个公式说明参数改变如何作用于面积密度；若还想保留 \(\det D\phi\) 的符号，就必须给参数域与曲面指定定向。

% !TeX root = ../../Book1.tex
\section{定向、体积形式与参数选择}
\label{b1:sec:B105}

面积公式中的 Jacobi 因子总是非负数，而通常行列式可以为负。负号并不是需要修正的面积误差；它记录了映射是否反转方向。对标量函数作曲面积分时方向可以忽略，对交错型积分或讨论曲面的边界时却不能。本节先在一个向量空间内固定符号规则，再将它用于连续变化的切空间。

\subsection{顶次外空间中的两个方向}
\label{b1:subsec:B10501}

\begin{definition}\label{b1:def:app-b-orientation}
设实向量空间 \(V\) 的维数为 \(n\geq1\)。对两组有序基 \(e=(e_1,\ldots,e_n)\) 与 \(f=(f_1,\ldots,f_n)\)，若由
\[
 f_j=\sum_{i=1}^n a_{ij}e_i
\]
定义的过渡矩阵满足 \(\det A>0\)，则称这两组基同向。选择全部有序基在此关系下的一个等价类，称为给 \(V\) 指定一个\term[dingxiang]{定向}[orientation]；被选类中的基称为正向基，另一类中的基称为负向基。
\end{definition}

行列式的乘法公式说明同向关系满足自反性、对称性与传递性。固定一组基后，其他基的过渡矩阵行列式只能为正或为负；交换固定基的前两个向量，或把第一个向量取负，可以互换这两种符号。因此 \(V\) 恰有两个定向。这个选择不依赖长度，也不需要内积。

\ProseParagraphBreak

外代数把这一定义压缩到一个一维空间中。由于
\[
 f_1\wedge\cdots\wedge f_n
  =(\det A)e_1\wedge\cdots\wedge e_n,
\]
同向的基恰好给出 \(\bigwedge^nV\setminus\{0\}\) 中同一条开射线。指定定向等于从这条实直线的两个方向中选一个方向，而不是选定某个特定长度的向量。

\ProseParagraphBreak

对任意线性自映射 \(A:V\rightarrow V\)，顶次外映射作用在一维空间上，必是乘以某个标量。基下的子式公式给出
\[
 \bigwedge\nolimits^n A=(\det A)I_{\bigwedge^nV}.
\]
因此行列式可以不借助具体矩阵而定义；换基后标量不变，乘法公式来自外映射与复合相容。若 \(A\) 可逆，正行列式对应保持定向，负行列式对应反转定向。奇异映射把顶次外空间压为零，此时不能谈它把一个定向送到了像的哪个 \(n\) 维定向。

\ProseParagraphBreak

给 \(V\) 增加欧氏内积后，每个定向都有唯一的正向单位顶次外向量
\[
 \tau_V=e_1\wedge\cdots\wedge e_n,
\]
其中 \(e_i\) 为任意正向标准正交基。另一组同向标准正交基的过渡矩阵行列式为 \(1\)，故 \(\tau_V\) 与基无关。反转定向只把 \(\tau_V\) 换成 \(-\tau_V\)。

\begin{proposition}\label{b1:prop:app-b-volume-form}
在已定向的 \(n\) 维欧氏空间 \(V\) 上，存在唯一的顶次交错型 \(\mathrm{vol}_V\)，使
\[
 \langle\mathrm{vol}_V,\tau_V\rangle=1.
\]
它称为与内积及定向相应的\term[tiji-xingshi]{体积形式}[volume form]。若 \(v_j=\sum_i a_{ij}e_i\)，其中 \(e_i\) 是正向标准正交基，则
\[
 \mathrm{vol}_V(v_1,\ldots,v_n)=\det A,\quad
 \bigl|\mathrm{vol}_V(v_1,\ldots,v_n)\bigr|
       =|v_1\wedge\cdots\wedge v_n|.
\]
若 \(V,W\) 都是已定向的 \(n\) 维欧氏空间，\(L:V\rightarrow W\) 线性，则
\[
 L^*\mathrm{vol}_W=(\det_{\mathrm{or}}L)\mathrm{vol}_V,
 \quad
 |\det_{\mathrm{or}}L|=J_n(L),
\]
其中 \(\det_{\mathrm{or}}L\) 是 \(L\) 在两端正向标准正交基下的行列式。
\end{proposition}
\begin{proof}
顶次外空间是一维的，所以在非零向量 \(\tau_V\) 上指定值便唯一决定一个线性泛函；再由外空间与交错型的配对识别得到 \(\mathrm{vol}_V\)。将
\(v_1\wedge\cdots\wedge v_n=(\det A)\tau_V\) 代入，就得到第一个公式。由于 \(|\tau_V|=1\)，取绝对值得到第二个公式。

最后，\(L^*\mathrm{vol}_W\) 仍是 \(V\) 上的顶次型，故为 \(\mathrm{vol}_V\) 的标量倍数。在正向标准正交基上求值，倍数恰为所写行列式。其绝对值由 Gram 公式等于面积因子。两端正向标准正交基的改变各自只有行列式 \(1\)，所以这个带符号的数不随所用基而变。
\end{proof}

欧氏体积测度与体积形式要区分。前者不给非负函数的积分增加符号；后者是作用于有序向量组的交错型，改变定向便取负。若把内积换掉而保持定向，正向基的集合不变，单位顶次外向量和体积形式的具体大小则会改变。

\ProseParagraphBreak

以下不是由前述有序基的等价类定义推出，而是本书为零维定向边界单独采用的符号约定：令其顶次外空间为 \(\bigwedge^0\{0\}\coloneqq\mathbb R\)，并用符号 \(+1\) 或 \(-1\) 指定定向。唯一的空基不能区分这两个选择。因而一个定向点带有正号或负号；有限多个定向点的积分就是带这些符号的求和。

\subsection{切空间的连续定向}
\label{b1:subsec:B10502}

设 \(M\subseteq\mathbb R^N\) 是 \(k\) 维 \(C^1\) 子流形。这里的局部描述是：每点附近都可由一个 \(C^1\) 参数化
\[
 f:U\subseteq\mathbb R^k\longrightarrow M
\]
给出，\(f\) 在该邻域是到其像的同胚，且 \(Df\) 处处具有秩 \(k\)。切空间为 \(T_{f(s)}M=Df(s)(\mathbb R^k)\)。本节只使用这类局部参数图及其 \(C^1\) 过渡映射；不在这里展开一般抽象流形理论。

\ProseParagraphBreak

参数域的标准定向给出局部连续的单位可分解外向量场
\[
 \tau_f\bigl(f(s)\bigr)
 =\frac{\partial_1f(s)\wedge\cdots\wedge\partial_kf(s)}
        {J_kf(s)}.
\]
分母处处为正，所以该式确实连续。若同一部分曲面另用 \(g=f\circ\phi\) 参数化，其中 \(\phi\) 为两个参数开集之间的 \(C^1\) 微分同胚，则
\begin{equation}\label{b1:eq:app-b-param-orientation}
 \tau_g\bigl(g(t)\bigr)
   =\operatorname{sgn}\bigl(\det D\phi(t)\bigr)
      \tau_f\bigl(g(t)\bigr).
\end{equation}
这是上一节外映射的复合公式与 \(J_k(g)=J_k(f\circ\phi)\) 的直接商式。过渡行列式连续且从不为零，所以其符号在每个连通分支上固定。

\ProseParagraphBreak

若能在 \(M\) 上连续选择单位可分解外向量 \(\tau(x)\)，使其所确定的平面为 \(T_xM\)，则称 \(M\) 可定向；选定这样的 \(\tau\) 就给出了 \(M\) 的定向。局部参数图若满足 \(\tau_f=\tau\)，便是正向参数图。由公式\eqref{b1:eq:app-b-param-orientation}，正向参数图之间的过渡行列式为正；反过来，若选出的参数图覆盖 \(M\)，且所有过渡映射的行列式均为正，那么这些局部 \(\tau_f\) 在重叠处相等，可以拼成全域连续的 \(\tau\)。

\ProseParagraphBreak

每张参数图都可定向，并不保证这些选择能在整张曲面上拼接。问题出在沿着不同图走一圈后符号能否保持一致，不出在某一点的线性代数。若 \(M\) 连通且已有一个定向 \(\tau\)，任何其他连续定向只能写成 \(\sigma(x)\tau(x)\)，其中 \(\sigma(x)\in\{-1,1\}\) 连续。连通性使 \(\sigma\) 为常数，故这时恰有两个全域定向。若有多个连通分支，各分支可以独立选择符号。

\ProseParagraphBreak

在\chapref{b1:ch:19}的可整流情形，切空间一般只在几乎处处存在；\appref{b1:app:D}将用可测的单位可分解外向量记录方向，而不要求它连续。连续定向与可测定向具有不同的用途，不能把流形的全局可定向性当作所有可整流积分对象的先决条件。

\subsection{外法向优先的边界约定}
\label{b1:subsec:B10503}

设 \(V\) 是已定向的 \(n\) 维欧氏空间，\(\nu\) 是一个单位向量，令 \(H=\nu^\perp\)。给 \(H\) 选择定向时，规定把 \(\nu\) 放在其正向基的最前面，所得到的 \(V\) 的基应为正向。这称为外法向优先约定；此处“外”是否具有几何含义，取决于 \(\nu\) 是否来自某个区域的外法向。

\begin{proposition}\label{b1:prop:app-b-boundary-orientation}
当 \(n\geq2\) 时，存在唯一的单位可分解外向量
\(\tau_H\in\bigwedge^{n-1}H\)，使
\[
 \nu\wedge\tau_H=\tau_V.
\]
它对应于上述定向。其体积形式由
\[
 \mathrm{vol}_H(v_1,\ldots,v_{n-1})
   =\mathrm{vol}_V(\nu,v_1,\ldots,v_{n-1}),
 \quad v_i\in H
\]
给出。反转 \(\nu\) 而保持 \(V\) 的定向，或反转 \(V\) 的定向而保持 \(\nu\)，都会反转 \(H\) 的定向。
\end{proposition}
\begin{proof}
选 \(H\) 的一组标准正交基 \(q_1,\ldots,q_{n-1}\)，则
\(\nu,q_1,\ldots,q_{n-1}\) 是 \(V\) 的标准正交基，故其外积为 \(\pm\tau_V\)。如有必要把 \(q_1\) 取负，就可令符号为正。由此构造了所需的 \(\tau_H\)。

空间 \(\bigwedge^{n-1}H\) 为一维，且映射 \(\eta\mapsto\nu\wedge\eta\) 把其非零基向量送到非零向量，因此单射，证明唯一性。右侧定义的型在这组正向标准正交切向量上取值为 \(1\)，所以按体积形式的唯一性就是 \(\mathrm{vol}_H\)。最后两种符号改变分别代入 \(\nu\wedge\tau_H=\tau_V\)，便得到相应的 \(\tau_H\) 都取负。
\end{proof}

若 \(\Omega\subseteq\mathbb R^n\) 有 \(C^1\) 边界，采用环境标准定向，并把连续单位外法向 \(\nu_\Omega\) 代入该命题，就得到边界定向。例如局部写
\[
 \Omega=\{x:F(x)<0\},\quad \nabla F\ne0\quad\text{于边界},
\]
则 \(\nu_\Omega=\nabla F/|\nabla F|\)。使用另一个表示同一内部一侧的定义函数，不会改变外法向。该约定与\chapref{b1:ch:27}的外法向一致；有限周长集合上的相应定向须在约化边界上作几乎处处的解释。

\ProseParagraphBreak

在平面单位圆上，外法向为
\(\nu(\theta)=(\cos\theta,\sin\theta)\)，而
\[
 \tau(\theta)=(-\sin\theta,\cos\theta)
\]
满足 \(\det(\nu,\tau)=1\)，所以正向是逆时针。若考虑一个圆环，内边界的区域外法向指向圆心，故内圆的正向是顺时针。两个边界分支的方向相反，来自同一个“区域的外侧”条件，不是额外补出的符号规则。

\ProseParagraphBreak

对图面 \(f(s)=(s,g(s))\subset\mathbb R^{k+1}\)，向上的单位法向为
\[
 \nu=\frac{(-\nabla g,1)}{\sqrt{1+|\nabla g|^2}}.
\]
列矩阵 \((\partial_1f,\ldots,\partial_kf,\nu)\) 的行列式为
\(\sqrt{1+|\nabla g|^2}>0\)。将最后一列移到第一列需要交换 \(k\) 次，因此外法向优先约定给出的切向外向量是
\[
 \tau_H=(-1)^k
 \frac{\partial_1f\wedge\cdots\wedge\partial_kf}
      {\sqrt{1+|\nabla g|^2}}.
\]
二维曲面的熟悉规则与一维曲线的规则因此看起来不同：关键是法向放在首位还是末位，不能在维数改变时省去这一次符号核对。

\ProseParagraphBreak

在 \(n=1\) 时把上述方程解释为 \(\nu\,\tau_H=\tau_V\)，其中 \(\tau_H\in\{-1,1\}\)。标准定向区间 \([a,b]\) 的左端外法向为 \(-1\)，右端为 \(1\)，所以其定向边界为“右端点减左端点”。这正与微积分基本定理中的 \(f(b)-f(a)\) 对应。

\subsection{定向积分与参数变换}
\label{b1:subsec:B10504}

现在设 \(\omega(x)\) 是定义在曲面附近的连续 \(k\) 次交错型场，即各点都有一个交错型，且其在固定坐标基上的系数连续。给曲面 \(M\) 指定单位切向外向量场 \(\tau\)，在
\[
 \int_M\bigl|\langle\omega(x),\tau(x)\rangle\bigr|
       \,\mathrm d\mathcal H^k(x)<\infty
\]
时，定义这个型场在定向曲面上的积分为
\[
 \int_{(M,\tau)}\omega
   \coloneqq\int_M\langle\omega(x),\tau(x)\rangle
          \,\mathrm d\mathcal H^k(x).
\]
右侧是本书已经建立的标量积分，所以定义先有明确的可积性条件，再讨论参数表达。这里不要求提前掌握微分形式的外微分或 Stokes 定理。

\ProseParagraphBreak

若 \(f:U\rightarrow M\) 是覆盖所考察曲面片的正向单射 \(C^1\) 参数图，则 \(f\) 在开集 \(U\) 上局部 Lipschitz。先把\cref{b1:cor:local-area-formulas}的非负加权局部面积公式用于
\(\bigl|\langle\omega(f(s)),\tau(f(s))\rangle\bigr|\)，由上面的可积性假设得到相应参数域积分有限；再使用同一结果的实值加权版本，便有
\[
 \begin{aligned}
 \int_{(M,\tau)}\omega
 &=\int_U
   \langle\omega\bigl(f(s)\bigr),\tau\bigl(f(s)\bigr)\rangle
       J_kf(s)\,\mathrm ds\\
 &=\int_U
   \omega\bigl(f(s)\bigr)(\partial_1f(s),\ldots,\partial_kf(s))
       \,\mathrm ds.
 \end{aligned}
\]
分母中的面积因子被抵消，剩下的是直接测试有序偏导向量。若改用负向参数图，同一个等式的最后一行要取负。

\ProseParagraphBreak

具体地，设 \(g=f\circ\phi\)，且 \(\det D\phi\) 在所用参数域上恒为负。型的多线性给出
\[
 \omega\bigl(g(t)\bigr)(\partial_1g(t),\ldots,\partial_kg(t))
 =(\det D\phi(t))\,
   \omega\Bigl(f\bigl(\phi(t)\bigr)\Bigr)
       \Bigl(\partial_1f\bigl(\phi(t)\bigr),\ldots,
             \partial_kf\bigl(\phi(t)\bigr)\Bigr).
\]
普通变量代换使用 \(|\det D\phi|\)，因而直接把 \(g\) 当作正向参数图来积分，就会得到原积分的负值。若保持曲面的同一个定向，则必须同时记入 \(g\) 的负向符号，两者抵消后积分不变。若有多个参数分支，分别在符号固定的分支上讨论即可。

\ProseParagraphBreak

这说明无向面积公式中的绝对值和交错型拉回中的带符号行列式并不矛盾。前者改变积分所用的非负测度，后者改变被积对象对有序向量的取值。将定向整体反转，面积测度不动而 \(\tau\) 取负，所以型的积分整体取负。\appref{b1:app:D}把这种积分看作作用于测试形式的线性泛函，边界也将在这一对偶语言中统一处理。


\begin{chaptersummary}
张量积通过商去双线性关系，把多线性问题转为线性问题；
外代数进一步商去重复因子，保留有序向量组的交错信息。
这两个构造都先于内积，不能把向量、对偶向量及坐标数组直接混为一谈。

\ProseParagraphBreak

欧氏内积传到外空间以后，单位正交外基的约定决定了 Gram 公式及体积归一化。
可分解外向量能够记录一个平面、一个体积尺度和一个方向，
一般外向量则未必对应单独一块平行多面体。
这种差别也解释了为何流的质量不能未经说明就换成系数的欧氏长度。

\ProseParagraphBreak

线性映射在各次外空间上的作用由子式组成，奇异值控制各维最大伸缩。
面积公式考察定义域的全部方向，余面积公式考察核的正交补；
在非正交坐标下，两端度量的 Gram 矩阵都必须纳入。
保留顶次外作用的符号便得到定向，取其长度才回到非负的 Jacobi 因子。
在有向积分中同时记录参数的符号和测度的绝对值，才能保持同一个几何对象不随表示改变。
\end{chaptersummary}

\BookExercises

% 习题均为自拟。每题具体训练增量见相邻注释；不将正文主构造留作习题。
\begin{exercise}\label{b1:ex:app-b-tensor-rank}
设 \(V,W\) 是有限维实空间，\(\xi\in V\otimes W\)。定义
\[
 T_\xi:V^*\longrightarrow W,\quad
 T_{v\otimes w}(\lambda)=\lambda(v)\cdot w,
\]
并按 \(\xi\) 线性延拓。证明该定义良定，且把 \(\xi\) 写成纯张量之和所需的最少项数等于 \(\operatorname{rank}T_\xi\)；零张量的最少项数规定为零。
对 \(V=\mathbb R^m,W=\mathbb R^n\)，把
\(\xi=\sum_{i,j}a_{ij}e_i\otimes f_j\) 的这一数值与矩阵 \((a_{ij})\) 的秩比较，并解释为何它不依赖所选基。
\end{exercise}
% 来源：自拟；由正文的非纯张量例子推进到最少分解项数的内在刻画。
\begin{solution}
对固定 \(\lambda\)，求值 \((v,w)\mapsto\lambda(v)\cdot w\) 双线性，因而可在张量积上线性化。
对每个 \(\xi\)，所得值又对 \(\lambda\) 线性，所以确实得到 \(T_\xi\)。
若 \(\xi=\sum_{\ell=1}^r v_\ell\otimes w_\ell\)，则
\[
 T_\xi(\lambda)=\sum_{\ell=1}^r\lambda(v_\ell)\cdot w_\ell.
\]
它的像包含在 \(\operatorname{span}(w_1,\ldots,w_r)\) 中，故
\(\operatorname{rank}T_\xi\leq r\)。

反过来，设像空间维数为 \(q\)，取其基 \(w_1,\ldots,w_q\)。
存在唯一的线性泛函 \(\ell_j\) 作用于 \(V^*\)，使
\[
 T_\xi(\lambda)=\sum_{j=1}^q\ell_j(\lambda)\cdot w_j.
\]
有限维双对偶的自然识别给出 \(v_j\in V\)，满足
\(\ell_j(\lambda)=\lambda(v_j)\)：若 \(e_i,e^i\) 互为对偶基，
可直接取 \(v_j=\sum_i\ell_j(e^i)\cdot e_i\) 并验证。
于是 \(\xi-\sum_jv_j\otimes w_j\) 对每个 \(e^i\) 的收缩都为零。
按基 \(e_i\otimes f_j\) 展开，这迫使全部系数为零，故得到了 \(q\) 项分解。
当 \(q=0\) 时同一坐标检验说明 \(\xi=0\)。

最后，\(T_\xi(e^i)=\sum_ja_{ij}f_j\)，所以 \(T_\xi\) 的矩阵是 \((a_{ij})\) 的转置。
两者秩相同，而线性映射像空间的维数不依赖基。
因此最少分解项数虽然可以通过一个系数矩阵计算，结果是张量本身的性质。
\end{solution}

\begin{exercise}\label{b1:ex:app-b-tensor-kernel}
设 \(A:V\rightarrow V'\)、\(B:W\rightarrow W'\) 线性。
把子空间的张量积通过自然映射识别为 \(V\otimes W\) 的子空间，证明
\[
 \ker(A\otimes B)
  =(\ker A)\otimes W+V\otimes(\ker B).
\]
进一步证明
\[
 \operatorname{rank}(A\otimes B)
   =(\operatorname{rank}A)(\operatorname{rank}B).
\]
说明为什么若只在纯张量上核验“一个因子落在核中”，还不足以证明第一个等式。
\end{exercise}
% 来源：自拟；将泛性质连接到核与商空间，并训练一般张量上的论证。
\begin{solution}
取直和分解
\[
 V=K\oplus V_0,\quad W=L\oplus W_0,\quad
 K=\ker A,\quad L=\ker B.
\]
分别选各分量的基并合在一起。由张量积的基定理，自然包含映射在对应基张量上单射，且
\[
 V\otimes W=(K\otimes L)\oplus(K\otimes W_0)
       \oplus(V_0\otimes L)\oplus(V_0\otimes W_0).
\]
前三项都被 \(A\otimes B\) 消去，其和恰为题述右侧。
最后一项中，\(A|_{V_0}\)、\(B|_{W_0}\) 分别是到各自像空间的线性同构。
它们的逆也诱导张量映射，并由复合公式给出逆映射，
故 \(A\otimes B\) 在 \(V_0\otimes W_0\) 上单射。
这同时说明没有其他核向量，并且其像为
\(\operatorname{ran}A\otimes\operatorname{ran}B\)，维数为两秩之积。

一般张量是若干纯张量的和，各项的像可能相互抵消。
对单项的判断没有排除这类抵消产生新的核元素；
上述直和分解将所有剩余项放在一个映射单射的子空间内，才完成了排除。
\end{solution}

\begin{exercise}\label{b1:ex:app-b-plucker-four}
在 \(\mathbb R^4\) 中写
\[
 \xi=\sum_{1\leq i<j\leq4}p_{ij}e_i\wedge e_j.
\]
证明 \(\xi\) 可分解当且仅当
\[
 p_{12}p_{34}-p_{13}p_{24}+p_{14}p_{23}=0.
\]
当 \(p_{12}\ne0\) 时给出显式的两个因子，并说明如何处理其余情形。
\end{exercise}
% 来源：自拟；将正文的不可分解反例推进到四维二次外向量的完整判据。
\begin{solution}
直接按交错法则展开得到
\[
 \xi\wedge\xi
 =2(p_{12}p_{34}-p_{13}p_{24}+p_{14}p_{23})
       e_1\wedge e_2\wedge e_3\wedge e_4.
\]
若 \(\xi=v\wedge w\)，则与自身的外积中有重复向量，故为零，证明必要性。

若 \(a=p_{12}\ne0\)，令
\[
 v=e_1-\frac{p_{23}}a e_3-\frac{p_{24}}a e_4,\quad
 w=ae_2+p_{13}e_3+p_{14}e_4.
\]
外积 \(v\wedge w\) 的前五个系数分别为
\(p_{12},p_{13},p_{14},p_{23},p_{24}\)，
而 \(e_3\wedge e_4\) 的系数为
\[
 \frac{p_{13}p_{24}-p_{14}p_{23}}a=p_{34}.
\]
因此 \(v\wedge w=\xi\)。
若 \(\xi\ne0\) 但 \(p_{12}=0\)，选取一个非零 \(p_{ij}\)，
将该对坐标排列到前两个位置。
条件 \(\xi\wedge\xi=0\) 与坐标排列无关，故在新基下可用刚才的构造，
再将所得向量写回原坐标。
若 \(\xi=0\)，取一个因子为零即可。这里四维二次的特殊形式很重要；
不能由这一题便断言所有次数与维数都有同样的一条二次判据。
\end{solution}

\begin{exercise}\label{b1:ex:app-b-projection-volume}
设 \(v_1,\ldots,v_k\in\mathbb R^n\)，\(1\leq k\leq n\)，
令 \(Q\) 为它们生成的平行多面体。
对大小为 \(k\) 的递增指标组 \(I\)，记 \(\pi_I\) 为到对应坐标 \(k\) 维平面的正交投影。
证明
\[
 \bigl(\mathcal H^k(Q)\bigr)^2
   =\sum_{|I|=k}\bigl(\mathcal H^k(\pi_IQ)\bigr)^2.
\]
在 \(n=3,k=2\) 时证明常数 \(1/\sqrt3\) 是不等式
\[
 \max_{|I|=2}\mathcal H^2(\pi_IQ)
   \geq\frac1{\sqrt3}\mathcal H^2(Q)
\]
中最好的统一常数，并构造一个非退化的取等例子。
\end{exercise}
% 来源：自拟；将子式系数解释为所有坐标投影的面积，并要求最优性构造。
\begin{solution}
记列坐标矩阵为 \(A\)。投影多面体的面积是
\(|\det A_I|\)，而 \(Q\) 的面积平方为
\(\det(A^{\mathsf T}A)\)。
外积子式展开和外基的正交性给出
\[
 |v_1\wedge\cdots\wedge v_k|^2=\sum_{|I|=k}|\det A_I|^2,
\]
所以得到所需几何恒等式。向量相关时，各项都为零，仍然成立。

当 \(n=3,k=2\) 时右侧有三项，最大一项的平方至少为总和的三分之一。
取
\[
 v_1=e_1-e_3,\quad v_2=e_2-e_3.
\]
它们的外积为
\(e_1\wedge e_2-e_1\wedge e_3+e_2\wedge e_3\)，
三个投影面积都是 \(1\)，原面积为 \(\sqrt3\)。
因此这一例子取等，任何更大的统一常数都会失败。
同一恒等式并不说原面积等于各投影面积之和；本题中后者是 \(3\)。
\end{solution}

\begin{exercise}\label{b1:ex:app-b-normal-transform}
设 \(A:\mathbb R^n\rightarrow\mathbb R^n\) 可逆，\(n\geq2\)。
给定单位向量 \(\nu\)，令 \(H=\nu^\perp\)，并按外法向优先约定赋予 \(H\) 定向 \(\tau_H\)。
证明像平面 \(A(H)\) 的单位法向可取
\[
 \nu'=\frac{A^{-\mathsf T}\nu}{|A^{-\mathsf T}\nu|}.
\]
用 \(\nu'\) 和环境标准定向确定 \(\tau_{A(H)}\)，证明
\[
 (\bigwedge\nolimits^{n-1}A)\tau_H
   =(\det A)|A^{-\mathsf T}\nu|\,\tau_{A(H)},
\]
并推出超平面上的面积伸缩公式
\[
 J_{n-1}(A;H)=|\det A|\cdot|A^{-\mathsf T}\nu|.
\]
\end{exercise}
% 来源：自拟；由顶次行列式推导余维一法向和面积的共同变换，不依赖未定义的余子式算子。
\begin{solution}
对 \(v\in H\)，有
\(\langle A^{-\mathsf T}\nu,Av\rangle=\langle\nu,v\rangle=0\)。
由于 \(A\) 可逆，\(A^{-\mathsf T}\nu\ne0\)，所以所写 \(\nu'\) 确为像平面的单位法向。

选 \(H\) 的正向标准正交基 \(q_1,\ldots,q_{n-1}\)。
由 \(\nu\wedge\tau_H=e_1\wedge\cdots\wedge e_n\)，
\((\nu,q_1,\ldots,q_{n-1})\) 为正向标准正交基。
像外向量落在一维空间 \(\bigwedge^{n-1}A(H)\) 中，故等于
\(c\tau_{A(H)}\)。把 \(\nu'\) 放在首位并取环境体积形式求值，有
\[
 \begin{aligned}
 c&=\det(\nu',Aq_1,\ldots,Aq_{n-1})\\
  &=(\det A)\det(A^{-1}\nu',q_1,\ldots,q_{n-1})\\
  &=(\det A)\langle A^{-1}\nu',\nu\rangle\\
  &=(\det A)|A^{-\mathsf T}\nu|.
 \end{aligned}
\]
第三行将第一列分为法向分量与切向分量，切向部分与其余列线性相关而消失。
取外向量长度，便得到面积伸缩公式。
特别当 \(\det A<0\) 时，直接传递原切向外向量与按像平面法向重新确定的定向相反；
只写非负面积公式会丢掉这一信息。
\end{solution}

\begin{exercise}\label{b1:ex:app-b-hodge-linear-map}
设 \(V\) 是已定向的 \(n\) 维欧氏空间，\(\tau_V\) 为其正向单位顶次外向量。
对 \(0\leq k\leq n\)，证明存在唯一线性映射
\[
 \star:\bigwedge\nolimits^kV\longrightarrow
             \bigwedge\nolimits^{n-k}V
\]
满足
\[
 \alpha\wedge(\star\beta)=\langle\alpha,\beta\rangle\tau_V
 \quad(\alpha,\beta\in\bigwedge\nolimits^kV).
\]
给出它在正向标准正交外基上的公式，并证明
\[
 |\star\beta|=|\beta|,\quad
 \star(\star\beta)=(-1)^{k(n-k)}\beta.
\]
最后判断反转 \(V\) 的定向时，这一映射如何变化。
\end{exercise}
% 来源：自拟；以非退化外配对构造 Hodge 型线性映射，增加对偶和定向的综合训练。
\begin{solution}
取正向标准正交基 \(e_1,\ldots,e_n\)。
对递增指标组 \(I\)，令 \(I^c\) 为递增补指标组，选
\(s_I\in\{-1,1\}\) 使
\[
 e_I\wedge(s_Ie_{I^c})=\tau_V.
\]
规定 \(\star e_I\coloneqq s_Ie_{I^c}\)，并线性延拓。
若 \(J\ne I\) 且 \(|J|=|I|=k\)，则 \(J\) 至少有一个指标属于 \(I^c\)，
所以 \(e_J\wedge e_{I^c}=0\)。
因此在任意两个基向量上都有
\(e_J\wedge\star e_I=\delta_{JI}\tau_V\)，双线性给出要求的等式。

另一方面，若 \(\eta\in\bigwedge^{n-k}V\) 满足每个 \(e_J\wedge\eta=0\)，
把 \(\eta\) 按补指标基展开，每次测试 \(e_J\) 都挑出对应的唯一系数。
于是 \(\eta=0\)，证明外配对非退化，也证明满足要求的 \(\star\) 唯一。
所定义的映射把标准正交基送到带符号的标准正交基，故保持长度。

由块交换法则，
\[
 e_{I^c}\wedge e_I=(-1)^{k(n-k)}e_I\wedge e_{I^c},
\]
所以 \(s_{I^c}=(-1)^{k(n-k)}s_I\)，从而
\(\star\star e_I=(-1)^{k(n-k)}e_I\)。
线性延拓得到一般情形；\(k=0,n\) 同样由空外积约定包含在内。
反转定向把右侧 \(\tau_V\) 取负，而内积不变，
因此由唯一性，新映射恰为原映射的负号。
\end{solution}

\begin{exercise}\label{b1:ex:app-b-mobius-orientation}
取 \(R>\varepsilon>0\)，考虑参数化
\[
 F(\theta,t)=
 \bigl((R+t\cdot \cos(\theta/2))\cdot \cos\theta,\,
       (R+t\cdot \cos(\theta/2))\cdot \sin\theta,\,
       t\cdot \sin(\theta/2)\bigr),
 \quad |t|<\varepsilon.
\]
先说明它在每点具有秩 \(2\)，且
\(F(\theta+2\pi,-t)=F(\theta,t)\)。
把这条缝合关系形成的曲面称为 Möbius 带，证明它没有连续定向。
同时解释这不妨碍每个足够小的参数片具有连续定向。
\end{exercise}
% 来源：自拟；显式跟踪接缝的符号，检验局部定向与全局连续定向的差别。
\begin{solution}
写柱坐标方向为
\[
 e_r=(\cos\theta,\sin\theta,0),\quad
 e_\theta=(-\sin\theta,\cos\theta,0),\quad e_z=(0,0,1).
\]
有
\[
 F_t=\cos(\theta/2)e_r+\sin(\theta/2)e_z,
\]
\[
 F_\theta=(R+t\cdot \cos(\theta/2))e_\theta
      -\frac t2\cdot \sin(\theta/2)e_r
      +\frac t2\cdot \cos(\theta/2)e_z.
\]
第一个向量长度为 \(1\) 且没有 \(e_\theta\) 分量；
第二个的 \(e_\theta\) 分量大于 \(R-\varepsilon>0\)。
所以它们线性无关。把 \(\theta+2\pi,-t\) 代入公式，
利用半角正弦、余弦同时变号，就得到所述缝合关系。
由于径向坐标始终为正，像点的极角决定 \(\theta\) 模 \(2\pi\)；
固定这一选择后，径向与竖直分量确定 \(t\)。
因此除此缝合外没有其他重合，局部确为光滑曲面。

假设存在连续单位切向外向量场 \(\tau\) 给曲面定向。
沿中央圆 \(t=0\) 定义
\[
 \eta(\theta)=
 \frac{F_\theta(\theta,0)\wedge F_t(\theta,0)}
      {|F_\theta(\theta,0)\wedge F_t(\theta,0)|}.
\]
它连续且处处为单位切向外向量，所以
\(\tau\bigl(F(\theta,0)\bigr)=\sigma(\theta)\cdot\eta(\theta)\)，
其中 \(\sigma:[0,2\pi]\rightarrow\{-1,1\}\) 连续。
区间连通，故 \(\sigma\) 为常数。
但 \(F(2\pi,0)=F(0,0)\)，
\(F_\theta(2\pi,0)=F_\theta(0,0)\)，而
\(F_t(2\pi,0)=-F_t(0,0)\)，因此
\(\eta(2\pi)=-\eta(0)\)。
这迫使同一点的 \(\tau\) 等于自己的负号，与单位长度矛盾。

在任何不绕过接缝的小参数片上，归一化的 \(F_\theta\wedge F_t\) 都给出连续定向。
失败的只是沿闭合路径将这些局部选择拼成同一个连续选择，
不是局部外积不存在，也不是面积为零。
\end{solution}

\begin{exercise}\label{b1:ex:app-b-rectangle-boundary}
设 \(Q=[a,b]\times[c,d]\)，\(a<b,c<d\)，
函数 \(P,S\) 在 \(Q\) 的某邻域上为 \(C^1\)。
对 \(v=(v_1,v_2)\in\mathbb R^2\)，令一型场
\[
 \omega(x,y)(v)=P(x,y)v_1+S(x,y)v_2.
\]
按区域外法向优先的约定给四条边定向，证明
\[
 \int_{\partial Q}\omega
   =\int_Q(\partial_xS-\partial_yP)\,\mathrm dx\,\mathrm dy.
\]
将 \(Q\) 分成有限个轴平行小矩形，说明分别应用此公式再相加时，
内部边为何完全抵消，并解释这一抵消为什么依赖方向而不只是边的集合相同。
\end{exercise}
% 来源：自拟；以已知的基本定理与 Fubini 推导定向边界积分，展示局部拼接的符号机制。
\begin{solution}
正向依次是底边向右、右边向上、顶边向左、左边向下。
按定向积分的参数表达，
\[
 \begin{aligned}
 \int_{\partial Q}\omega
  &=\int_a^bP(x,c)\,\mathrm dx
     +\int_c^dS(b,y)\,\mathrm dy\\
  &\quad-\int_a^bP(x,d)\,\mathrm dx
     -\int_c^dS(a,y)\,\mathrm dy.
 \end{aligned}
\]
由一维微积分基本定理，
\[
 P(x,c)-P(x,d)=-\int_c^d\partial_yP(x,y)\,\mathrm dy,
\]
\[
 S(b,y)-S(a,y)=\int_a^b\partial_xS(x,y)\,\mathrm dx.
\]
导数在紧矩形上连续，因而可积；Fubini 定理允许把两个迭代积分写成同一个二重积分，
即得结论。四个顶点的长度测度为零，不影响边积分。

对相邻小矩形，它们在公共边上的外法向相反，
由边界定向约定，相应切向也相反。
沿同一几何线段的两个积分于是互为负数，逐对抵消，
仅留下原矩形的外边界。右侧则按不重叠内部作积分相加，得到整个 \(Q\) 的积分。
如果把边积分误写成相对于无向长度的同一个正密度积分，
内部边会相加成两倍而非抵消。边界作为集合没有记录足够的信息；
交错型与定向共同提供了这里所需的符号。
\end{solution}
